\chapter{Graphs: Shortest Path Algorithms}
\chapterepigraph{Plain BFS already computes shortest paths -- but only
when every edge costs exactly the same. The moment edges carry
different weights, BFS's level-by-level guarantee breaks, and a new
idea is needed: always expand whichever frontier node currently has
the smallest known distance. That single idea, plus two ways of
handling weights BFS and that idea still can't manage (negative edges,
and all pairs at once), is this entire chapter.}

\section{Relaxation: The One Operation Every Algorithm Here Repeats}
\label{sec:graph-shortest-path-intro}

Every algorithm in this chapter, however different its outer loop,
performs the same inner operation, called \textbf{relaxation}: for an
edge $u \to v$ with weight $w$, if
$\textit{dist}[u] + w < \textit{dist}[v]$, update
$\textit{dist}[v] \gets \textit{dist}[u] + w$ (a shorter path to $v$
was just found, via $u$). The algorithms differ only in \textbf{which
edges they relax, and in what order}.

\begin{patternbox}[What This Chapter Covers, and How Each Piece Differs]
\begin{itemize}
  \item \textbf{Shortest Path, Unit Weights (Section~\ref{sec:graph-shortest-path-unit-weights})}:
        plain BFS already \emph{is} correct relaxation when every edge
        weighs $1$ -- level order and shortest distance coincide
        exactly.
  \item \textbf{Dijkstra's Algorithm (Section~\ref{sec:graph-dijkstra})}:
        for non-negative weights, always relax edges out of whichever
        unfinished node currently has the \textbf{smallest known}
        distance, using a priority queue to find that node quickly.
  \item \textbf{Number of Ways to Arrive (Section~\ref{sec:graph-number-of-ways-arrive})}
        and \textbf{Path With Minimum Effort (Section~\ref{sec:graph-path-min-effort})}:
        Dijkstra's frontier-expansion order, reused with the
        relaxation rule itself changed -- counting ties instead of
        only keeping the minimum, and minimizing the \emph{maximum}
        edge on a path instead of the sum.
  \item \textbf{Cheapest Flights Within K Stops (Section~\ref{sec:graph-cheapest-flights-k-stops})}:
        relaxation bounded to a fixed number of rounds, since Dijkstra's
        ``always pick the smallest'' greedy choice can skip a path that
        is cheaper only because it uses more stops.
  \item \textbf{Bellman-Ford (Section~\ref{sec:graph-bellman-ford})}:
        relax \textbf{every} edge, $V{-}1$ times over, which is what
        makes it correct even with \textbf{negative} edge weights
        (where Dijkstra's greedy frontier choice can be wrong) -- and a
        $V$-th round of relaxation that still finds an improvement
        reveals a negative cycle.
  \item \textbf{Floyd-Warshall (Section~\ref{sec:graph-floyd-warshall})}
        and \textbf{City with Smallest Number of Neighbors (Section~\ref{sec:graph-city-smallest-neighbors})}:
        relax every pair of nodes through every possible intermediate
        node, computing \textbf{all-pairs} shortest distances at once.
\end{itemize}
\end{patternbox}

\newpage

\section{Shortest Path in an Undirected Graph with Unit Weights}
\label{sec:graph-shortest-path-unit-weights}

\problemstatement{Given an undirected graph where every edge has
weight $1$, and a source node, find the shortest distance from the
source to every other node.}

\begin{patternbox}
\emph{``Shortest path, every edge the same cost''} is exactly what
plain BFS already computes: BFS processes nodes in strictly
non-decreasing order of distance from the source, so the level at which
a node is first reached \textbf{is} its shortest distance. No
relaxation, priority queue, or weight comparison is needed at all.
\end{patternbox}

\subsection*{Why unit weights make BFS sufficient}

The general shortest-path relaxation rule,
$\textit{dist}[v] = \min(\textit{dist}[v], \textit{dist}[u]+w)$,
collapses to $\textit{dist}[v] = \textit{dist}[u]+1$ when every $w=1$ --
and BFS's queue already guarantees every node is dequeued with its
\emph{final}, minimal distance already correctly assigned (this is the
same property exploited by multi-source BFS in the Graph Basics
chapter's 01 Matrix problem, here with a single source instead of
many).

\subsection*{Algorithm}

\begin{enumerate}
  \item Initialize $\textit{dist}[\textit{src}]=0$ and every other
        $\textit{dist}[i]=\infty$.
  \item BFS from the source: for each neighbor of the current node
        with $\textit{dist}[\textit{neighbor}] > \textit{dist}[\textit{node}]+1$,
        update it and push the neighbor.
\end{enumerate}

\subsection*{C++ implementation}

\begin{lstlisting}
#include <bits/stdc++.h>
using namespace std;

vector<int> shortestPathUnitWeights(int n, vector<vector<int>>& adj, int src) {
    vector<int> dist(n, INT_MAX);
    dist[src] = 0;
    queue<int> q;
    q.push(src);

    while (!q.empty()) {
        int node = q.front();
        q.pop();

        for (int neighbor : adj[node]) {
            if (dist[node] + 1 < dist[neighbor]) {
                dist[neighbor] = dist[node] + 1;
                q.push(neighbor);
            }
        }
    }
    return dist;
}

int main() {
    int n = 5;
    vector<vector<int>> adj(n);
    auto addEdge = [&](int u, int v) { adj[u].push_back(v); adj[v].push_back(u); };
    addEdge(0, 1); addEdge(0, 2); addEdge(1, 3); addEdge(2, 3); addEdge(3, 4);

    vector<int> dist = shortestPathUnitWeights(n, adj, 0);
    for (int d : dist) cout << d << " ";
    cout << endl; // 0 1 1 2 3
    return 0;
}
\end{lstlisting}

\subsection*{Dry run}

\dryrun{Take the graph above, source $0$.}

$\textit{dist}=[0,\infty,\infty,\infty,\infty]$. Pop $0$: relax $1$
($\infty\to1$) and $2$ ($\infty\to1$); push both. Pop $1$: relax $3$
($\infty\to2$); push. Pop $2$: neighbor $3$ already at distance $2$,
and $\textit{dist}[2]+1=2$ is not strictly smaller, no update. Pop
$3$: relax $4$ ($\infty\to3$); push. Pop $4$: no unrelaxed neighbors.
Final: $[0,1,1,2,3]$.

\begin{complexitybox}
\textbf{Time Complexity.} $O(V+E)$.
\textbf{Space Complexity.} $O(V)$ for \textit{dist} and the queue.
\end{complexitybox}

\begin{takeawaybox}
This problem is included specifically as a baseline: every other
algorithm in this chapter is, at its core, solving the same
``relax edges until distances stabilize'' problem that unit-weight BFS
solves trivially -- they exist only because weighted, negative, or
all-pairs distances break BFS's clean level-by-level guarantee in
different specific ways.
\end{takeawaybox}

\section{Shortest Path in a Directed Acyclic Graph}
\label{sec:graph-shortest-path-dag}

\problemstatement{Given a weighted \textbf{directed acyclic graph}
(DAG) and a source node, find the shortest distance from the source to
every other node. Weights may even be negative -- the absence of cycles
is what makes this safe.}

\begin{patternbox}
\emph{``Shortest path'' + ``directed acyclic''} unlocks something
neither BFS nor Dijkstra needs: a \textbf{topological order}. If we
relax every node's outgoing edges in topological order, then by the
time we process a node its own shortest distance is already final,
because every path into it comes from nodes appearing \emph{earlier}
in the order.
\end{patternbox}

\subsection*{Why a topological order makes one pass enough}

In a DAG every edge $u \to v$ points from an earlier node to a later
node in topological order. So when we reach $u$ in that order, all
edges that could ever improve $\textit{dist}[u]$ have already been
relaxed -- they all originate from nodes before $u$. Thus
$\textit{dist}[u]$ is final the moment we start processing $u$, and a
single left-to-right sweep relaxing each node's out-edges suffices. No
priority queue and no repeated passes are needed, which is why this
beats both Dijkstra's ($O((V+E)\log V)$) and Bellman-Ford's
($O(VE)$) for DAGs, and tolerates negative weights that Dijkstra's
cannot.

\subsection*{Algorithm}

\begin{enumerate}
  \item Topologically sort the DAG (DFS-based, pushing nodes onto a
        stack after visiting all descendants).
  \item Initialize $\textit{dist}[\textit{src}]=0$, all others
        $\infty$.
  \item Pop nodes in topological order; for each node with a finite
        distance, relax every outgoing edge $u\to v$ of weight $w$: if
        $\textit{dist}[u]+w<\textit{dist}[v]$, update
        $\textit{dist}[v]$.
\end{enumerate}

\begin{ideabox}[Topological Order Replaces the Priority Queue]
Dijkstra's uses a heap to always pick the next node whose distance is
final. A DAG hands you that order for free: any topological order is
already a valid ``finalization order.'' Sorting once ($O(V+E)$) removes
the need to repeatedly find the minimum.
\end{ideabox}

\subsection*{C++ implementation}

\begin{lstlisting}
#include <bits/stdc++.h>
using namespace std;

void topoDFS(int node, vector<vector<pair<int,int>>>& adj,
             vector<bool>& vis, stack<int>& st) {
    vis[node] = true;
    for (auto& [nbr, w] : adj[node])
        if (!vis[nbr]) topoDFS(nbr, adj, vis, st);
    st.push(node);                       // push after all descendants
}

vector<int> shortestPathDAG(int n, vector<vector<pair<int,int>>>& adj, int src) {
    vector<bool> vis(n, false);
    stack<int> st;
    for (int i = 0; i < n; i++) if (!vis[i]) topoDFS(i, adj, vis, st);

    vector<int> dist(n, INT_MAX);
    dist[src] = 0;
    while (!st.empty()) {
        int node = st.top(); st.pop();   // topological order
        if (dist[node] == INT_MAX) continue;
        for (auto& [nbr, w] : adj[node])
            if (dist[node] + w < dist[nbr])
                dist[nbr] = dist[node] + w;   // relax
    }
    return dist;
}

int main() {
    int n = 6;
    vector<vector<pair<int,int>>> adj(n);   // directed {neighbor, weight}
    adj[0].push_back({1, 2}); adj[0].push_back({4, 1});
    adj[1].push_back({2, 3}); adj[2].push_back({3, 6});
    adj[4].push_back({2, 2}); adj[4].push_back({5, 4});
    adj[5].push_back({3, 1});

    vector<int> dist = shortestPathDAG(n, adj, 0);
    for (int d : dist) cout << (d == INT_MAX ? -1 : d) << " ";
    cout << endl;                        // 0 2 3 8 1 5
    return 0;
}
\end{lstlisting}

\subsection*{Dry run}

\dryrun{Source $0$; a topological order is $0,4,5,1,2,3$.}

\begin{itemize}
  \item $0$ ($\textit{dist}=0$): relax $1\to2$, $4\to1$.
  \item $4$ ($\textit{dist}=1$): relax $2\to1+2=3$, $5\to1+4=5$.
  \item $5$ ($\textit{dist}=5$): relax $3\to5+1=6$.
  \item $1$ ($\textit{dist}=2$): relax $2$ via $1$: $2+3=5>3$, no
        change.
  \item $2$ ($\textit{dist}=3$): relax $3$ via $2$: $3+6=9>6$, no
        change.
  \item $3$ ($\textit{dist}=6$): no out-edges. Final:
        $[0,2,3,6,1,5]$.
\end{itemize}

\subsection*{Complexity}

\begin{complexitybox}
\textbf{Time Complexity.} $O(V+E)$ -- one topological sort plus one
relaxation sweep, each linear. \textbf{Space Complexity.} $O(V)$ for
the distance array, visited array, and stack.
\end{complexitybox}

\begin{takeawaybox}
On a DAG, a topological order \emph{is} the order in which distances
become final, so a single linear sweep beats Dijkstra's and Bellman-Ford
and even handles negative weights. Whenever a shortest-path problem
promises ``no cycles,'' reach for topological sort before any heap.
\end{takeawaybox}

\section{Dijkstra's Algorithm}
\label{sec:graph-dijkstra}

\problemstatement{Given a weighted graph with \textbf{non-negative}
edge weights and a source node, find the shortest distance from the
source to every other node.}

\begin{patternbox}
\emph{``Shortest path, non-negative weights''} is solved by always
expanding the \textbf{unfinished} node with the currently
\textbf{smallest known} distance -- a greedy choice, made efficient by
storing candidate $(\textit{distance}, \textit{node})$ pairs in a
\textbf{min-heap} (priority queue) so the smallest is always available
in $O(\log V)$.
\end{patternbox}

\subsection*{Why the greedy choice is correct (and needs non-negative weights)}

When the node $u$ with the smallest tentative distance is popped, that
distance is guaranteed \textbf{final} -- no future relaxation through
any other, not-yet-finalized node $x$ could ever produce a shorter path
to $u$, because every such path would have to pass through some node
with tentative distance $\ge \textit{dist}[u]$ \emph{and then add a
non-negative edge weight}, which can only make the total $\ge
\textit{dist}[u]$. This argument breaks immediately if a negative edge
weight is allowed: a long detour through a very negative edge could
undercut the greedily-chosen shortest path, which is exactly why
Dijkstra's algorithm requires non-negative weights and Bellman-Ford
(Section~\ref{sec:graph-bellman-ford}) exists for the negative case.

\subsection*{Algorithm}

\begin{enumerate}
  \item Initialize $\textit{dist}[\textit{src}]=0$, all others
        $\infty$. Push $(0, \textit{src})$ into a min-heap.
  \item While the heap is non-empty: pop the smallest
        $(\textit{d}, \textit{node})$; if $\textit{d} >
        \textit{dist}[\textit{node}]$, this is a stale entry (a better
        path was already found), skip it.
  \item Otherwise, relax every edge out of $\textit{node}$: if
        $\textit{dist}[\textit{node}] + w < \textit{dist}[\textit{neighbor}]$,
        update and push $(\textit{dist}[\textit{neighbor}],
        \textit{neighbor})$.
\end{enumerate}

\begin{ideabox}[Why Stale Entries Are Allowed Rather Than Removed]
A node may be pushed onto the heap multiple times, once per
improvement to its tentative distance, rather than removing and
re-inserting it (which most heap implementations don't support
efficiently). The $\textit{d} > \textit{dist}[\textit{node}]$ check
simply ignores outdated copies when they're eventually popped --
correct, and simpler than maintaining a decrease-key operation.
\end{ideabox}

\subsection*{C++ implementation}

\begin{lstlisting}
#include <bits/stdc++.h>
using namespace std;

vector<int> dijkstra(int n, vector<vector<pair<int,int>>>& adj, int src) {
    vector<int> dist(n, INT_MAX);
    dist[src] = 0;
    priority_queue<pair<int,int>, vector<pair<int,int>>, greater<>> pq;
    pq.push({0, src});

    while (!pq.empty()) {
        auto [d, node] = pq.top();
        pq.pop();
        if (d > dist[node]) continue; // stale entry

        for (auto& [neighbor, weight] : adj[node]) {
            if (dist[node] + weight < dist[neighbor]) {
                dist[neighbor] = dist[node] + weight;
                pq.push({dist[neighbor], neighbor});
            }
        }
    }
    return dist;
}

int main() {
    int n = 5;
    vector<vector<pair<int,int>>> adj(n); // {neighbor, weight}
    auto addEdge = [&](int u, int v, int w) {
        adj[u].push_back({v, w});
        adj[v].push_back({u, w});
    };
    addEdge(0, 1, 4); addEdge(0, 2, 1); addEdge(2, 1, 2);
    addEdge(1, 3, 1); addEdge(2, 3, 5); addEdge(3, 4, 3);

    vector<int> dist = dijkstra(n, adj, 0);
    for (int d : dist) cout << d << " ";
    cout << endl; // 0 3 1 4 7
    return 0;
}
\end{lstlisting}

\subsection*{Dry run}

\dryrun{Take the graph above, source $0$.}

Heap: $[(0,0)]$. Pop $(0,0)$: relax $1$ ($\infty\to4$), $2$
($\infty\to1$); push both. Heap: $[(1,2),(4,1)]$. Pop $(1,2)$
(smallest): relax $1$ via $2$ ($4\to1+2=3$, improved); relax $3$ via
$2$ ($\infty\to1+5=6$); push $(3,1)$ and $(6,3)$. Heap:
$[(3,1),(4,1)_{\text{stale}},(6,3)]$. Pop $(3,1)$: relax $3$ via $1$
($6\to3+1=4$, improved); push $(4,3)$. Pop $(4,1)_{\text{stale}}$:
$4 > \textit{dist}[1]=3$, skip. Pop $(4,3)$: relax $4$ via $3$
($\infty\to4+3=7$); push. Pop $(6,3)_{\text{stale}}$: $6 >
\textit{dist}[3]=4$, skip. Pop $(7,4)$: no further relaxation. Final:
$\textit{dist}=[0,3,1,4,7]$.

\begin{complexitybox}
\textbf{Time Complexity.} $O((V+E)\log V)$ -- each edge can trigger one
heap push, and each heap operation costs $O(\log V)$.
\textbf{Space Complexity.} $O(V)$ for \textit{dist}, plus $O(E)$ for
the heap in the worst case (allowing stale duplicate entries).
\end{complexitybox}

\begin{takeawaybox}
Dijkstra's algorithm is BFS's natural generalization to weighted
graphs: BFS implicitly always expands the nearest unvisited node
because its queue processes things in arrival order, which only equals
distance order when every edge costs the same; Dijkstra's makes that
``always expand the nearest'' rule explicit via a priority queue,
restoring correctness for arbitrary non-negative weights.
\end{takeawaybox}

\section{Dijkstra's Algorithm Using a Set}
\label{sec:graph-dijkstra-set}

\problemstatement{Solve the same non-negative-weight shortest-path
problem as Section~\ref{sec:graph-dijkstra}, but use a \textbf{balanced
BST set} instead of a priority queue -- a variant that lets you actively
remove stale entries.}

\begin{patternbox}
The priority-queue Dijkstra tolerates stale duplicate entries and skips
them on pop. A \texttt{std::set} instead supports \textbf{erase}, so
when a node's distance improves we can \emph{delete} its old
$(\textit{dist},\textit{node})$ pair before inserting the new one --
emulating a decrease-key and keeping the structure free of outdated
copies.
\end{patternbox}

\subsection*{Set vs priority queue: the same algorithm, a different container}

The logic is identical to heap-based Dijkstra: repeatedly take the
smallest tentative-distance node, finalize it, and relax its edges. The
only difference is bookkeeping. A \texttt{set<pair<int,int>>} keeps
$(\textit{distance},\textit{node})$ pairs sorted, so
\texttt{begin()} is the minimum. When we find a shorter path to a
neighbor that already sits in the set with an older, larger distance, we
first \texttt{erase} that stale pair and then insert the improved one.
This keeps the set size bounded by $V$ rather than growing with every
improvement.

\begin{ideabox}[Erase-Then-Insert Emulates Decrease-Key]
The heap version cannot cheaply delete an interior element, so it lets
stale entries pile up and skips them later. The set version can delete
in $O(\log V)$, so it removes the obsolete pair up front. Both are
correct; the set trades a slightly larger constant factor for a smaller,
cleaner frontier.
\end{ideabox}

\subsection*{Algorithm}

\begin{enumerate}
  \item Initialize $\textit{dist}[\textit{src}]=0$, others $\infty$.
        Insert $(0,\textit{src})$ into the set.
  \item While the set is non-empty: take and erase the smallest pair
        $(\textit{d},\textit{node})$ from \texttt{begin()}.
  \item Relax each edge $\textit{node}\to\textit{nbr}$ of weight $w$: if
        $\textit{dist}[\textit{node}]+w<\textit{dist}[\textit{nbr}]$,
        erase the old $(\textit{dist}[\textit{nbr}],\textit{nbr})$ if
        present, update the distance, and insert the new pair.
\end{enumerate}

\subsection*{C++ implementation}

\begin{lstlisting}
#include <bits/stdc++.h>
using namespace std;

vector<int> dijkstraSet(int n, vector<vector<pair<int,int>>>& adj, int src) {
    vector<int> dist(n, INT_MAX);
    dist[src] = 0;
    set<pair<int,int>> st;               // {distance, node}, sorted
    st.insert({0, src});

    while (!st.empty()) {
        auto [d, node] = *st.begin();
        st.erase(st.begin());            // take the minimum

        for (auto& [nbr, w] : adj[node]) {
            if (dist[node] + w < dist[nbr]) {
                if (dist[nbr] != INT_MAX)
                    st.erase({dist[nbr], nbr});   // remove stale pair
                dist[nbr] = dist[node] + w;
                st.insert({dist[nbr], nbr});
            }
        }
    }
    return dist;
}

int main() {
    int n = 5;
    vector<vector<pair<int,int>>> adj(n);   // {neighbor, weight}
    auto addEdge = [&](int u, int v, int w) {
        adj[u].push_back({v, w});
        adj[v].push_back({u, w});
    };
    addEdge(0, 1, 4); addEdge(0, 2, 1); addEdge(2, 1, 2);
    addEdge(1, 3, 1); addEdge(2, 3, 5); addEdge(3, 4, 3);

    vector<int> dist = dijkstraSet(n, adj, 0);
    for (int d : dist) cout << d << " ";
    cout << endl;                        // 0 3 1 4 7
    return 0;
}
\end{lstlisting}

\subsection*{Dry run}

\dryrun{Same graph as Section~\ref{sec:graph-dijkstra}, source $0$.}

\begin{itemize}
  \item Set $\{(0,0)\}$. Take $(0,0)$; relax $1\!\to\!4$, $2\!\to\!1$.
        Set $\{(1,2),(4,1)\}$.
  \item Take $(1,2)$; relax $1$ via $2$: $3<4$, so erase $(4,1)$,
        insert $(3,1)$; relax $3\!\to\!6$. Set $\{(3,1),(6,3)\}$.
  \item Take $(3,1)$; relax $3$ via $1$: $4<6$, erase $(6,3)$, insert
        $(4,3)$. Set $\{(4,3)\}$.
  \item Take $(4,3)$; relax $4\!\to\!7$. Take $(7,4)$: nothing.
        Final $[0,3,1,4,7]$ -- no stale entries were ever skipped.
\end{itemize}

\subsection*{Complexity}

\begin{complexitybox}
\textbf{Time Complexity.} $O((V+E)\log V)$ -- same asymptotics as the
heap version; each edge may cause an erase and an insert, each
$O(\log V)$. \textbf{Space Complexity.} $O(V)$ -- the set never holds
more than one entry per node.
\end{complexitybox}

\begin{takeawaybox}
The set-based Dijkstra is the same greedy algorithm with a container
that supports deletion, so the frontier stays free of stale entries at
the cost of a larger constant. Use the priority-queue form by default;
reach for the set form when you specifically want an explicit
decrease-key or a frontier bounded by $V$.
\end{takeawaybox}

\section{Print the Shortest Path}
\label{sec:graph-print-shortest-path}

\problemstatement{Given a weighted undirected graph with non-negative
weights, find not just the shortest distance from node $1$ to node $n$,
but the actual \textbf{sequence of nodes} on one shortest path.}

\begin{patternbox}
\emph{``Print/return the path,'' not just its length}, means: run
Dijkstra's as usual, but additionally remember, for each node, the
\textbf{predecessor} through which its best distance was achieved.
Following those predecessor links backward from the destination
reconstructs the path.
\end{patternbox}

\subsection*{Recording where each improvement came from}

Every time relaxation improves $\textit{dist}[v]$ using edge $u\to v$,
node $u$ is the current best ``parent'' of $v$ on the shortest path.
Store $\textit{parent}[v]=u$ at that moment. When the algorithm
finishes, $\textit{parent}$ encodes a shortest-path tree rooted at the
source. Starting from the destination and repeatedly moving to the
parent walks the path in reverse; reversing that list gives the path
from source to destination.

\begin{ideabox}[A Parent Pointer Turns Distances into Paths]
Any shortest-\emph{distance} algorithm becomes a shortest-\emph{path}
algorithm by adding one array. The parent is updated in lockstep with
the distance -- exactly when, and only when, a strictly better distance
is found -- so it always reflects the edge on the current best path.
\end{ideabox}

\subsection*{Algorithm}

\begin{enumerate}
  \item Run Dijkstra's from the source. Initialize
        $\textit{parent}[i]=i$ for all $i$.
  \item On each successful relaxation of edge $u\to v$, set
        $\textit{parent}[v]=u$.
  \item If $\textit{dist}[\textit{dest}]=\infty$, no path exists.
        Otherwise start at \textit{dest}, follow \textit{parent} until
        reaching the source, collecting nodes, then reverse.
\end{enumerate}

\subsection*{C++ implementation}

\begin{lstlisting}
#include <bits/stdc++.h>
using namespace std;

vector<int> shortestPath(int n, vector<vector<pair<int,int>>>& adj, int src, int dest) {
    vector<int> dist(n + 1, INT_MAX), parent(n + 1);
    for (int i = 1; i <= n; i++) parent[i] = i;
    dist[src] = 0;
    priority_queue<pair<int,int>, vector<pair<int,int>>, greater<>> pq;
    pq.push({0, src});

    while (!pq.empty()) {
        auto [d, node] = pq.top(); pq.pop();
        if (d > dist[node]) continue;
        for (auto& [nbr, w] : adj[node]) {
            if (dist[node] + w < dist[nbr]) {
                dist[nbr] = dist[node] + w;
                parent[nbr] = node;          // remember predecessor
                pq.push({dist[nbr], nbr});
            }
        }
    }

    if (dist[dest] == INT_MAX) return {};    // unreachable
    vector<int> path;
    for (int cur = dest; cur != src; cur = parent[cur]) path.push_back(cur);
    path.push_back(src);
    reverse(path.begin(), path.end());
    return path;
}

int main() {
    int n = 5;
    vector<vector<pair<int,int>>> adj(n + 1);   // 1-indexed
    auto addEdge = [&](int u, int v, int w) {
        adj[u].push_back({v, w});
        adj[v].push_back({u, w});
    };
    addEdge(1, 2, 2); addEdge(1, 3, 4); addEdge(2, 3, 1);
    addEdge(2, 4, 7); addEdge(3, 5, 3); addEdge(5, 4, 2);

    vector<int> path = shortestPath(n, adj, 1, 4);
    for (int node : path) cout << node << " ";
    cout << endl;                            // 1 2 3 5 4
    return 0;
}
\end{lstlisting}

\subsection*{Dry run}

\dryrun{Source $1$, destination $4$.}

\begin{itemize}
  \item Dijkstra finalizes: $\textit{dist}[1]=0$, $[2]=2$ (parent
        $1$), $[3]=3$ (parent $2$, via $2\!+\!1$), $[5]=6$ (parent
        $3$), $[4]=8$ (parent $5$, via $6\!+\!2$, better than
        $2\!+\!7=9$).
  \item Backtrack from $4$: $4\to5\to3\to2\to1$.
  \item Reverse: $1,2,3,5,4$ (total distance $8$).
\end{itemize}

\subsection*{Complexity}

\begin{complexitybox}
\textbf{Time Complexity.} $O((V+E)\log V)$ -- Dijkstra's cost; the
backtrack is $O(V)$. \textbf{Space Complexity.} $O(V)$ for the distance
and parent arrays.
\end{complexitybox}

\begin{takeawaybox}
To turn any relaxation-based shortest-distance algorithm into one that
prints the path, record a parent pointer at each successful relaxation
and walk it backward from the destination. The same trick works for
BFS, Dijkstra's, and Bellman-Ford alike.
\end{takeawaybox}

\section{Shortest Distance in a Binary Maze}
\label{sec:graph-shortest-binary-maze}

\problemstatement{Given an $n\times m$ grid of $0$s (blocked) and $1$s
(open), a source cell and a destination cell, find the length of the
shortest path between them moving only through open cells in the four
cardinal directions. Return $-1$ if unreachable.}

\begin{patternbox}
A grid where every move costs the same one step is a
\textbf{unit-weight} graph, so plain \textbf{BFS} finds the shortest
path -- exactly the case from
Section~\ref{sec:graph-shortest-path-unit-weights}, just with the
graph given implicitly as a grid rather than an adjacency list. No
Dijkstra's is needed when all moves cost $1$.
\end{patternbox}

\subsection*{The grid is a graph in disguise}

Each open cell is a node; each pair of side-adjacent open cells is an
edge of weight $1$. Because BFS explores strictly in order of increasing
number of steps, the first time it reaches the destination it has done
so by a shortest path. We store $(\textit{row},\textit{col},
\textit{distance})$ in the queue (or track distance in a separate grid)
and expand the four neighbors, skipping blocked or out-of-bounds cells
and cells already visited.

\begin{ideabox}[BFS Layers Are Distance Layers]
Since every edge weighs $1$, BFS's $k$-th layer holds exactly the cells
at distance $k$ from the source. Marking a cell visited the moment it is
enqueued guarantees each cell is reached by its shortest distance and
never reprocessed.
\end{ideabox}

\subsection*{Algorithm}

\begin{enumerate}
  \item If the source is blocked, return $-1$ (or $0$ if source equals
        destination). Push the source with distance $0$; mark it
        visited.
  \item While the queue is non-empty, pop $(r,c,d)$; if it is the
        destination, return $d$.
  \item For each of the four neighbors that is in bounds, open, and
        unvisited: mark visited and push it with distance $d+1$.
  \item If the queue empties without reaching the destination, return
        $-1$.
\end{enumerate}

\subsection*{C++ implementation}

\begin{lstlisting}
#include <bits/stdc++.h>
using namespace std;

int shortestPathBinaryMaze(vector<vector<int>>& grid,
                           pair<int,int> src, pair<int,int> dest) {
    int n = grid.size(), m = grid[0].size();
    if (grid[src.first][src.second] == 0) return -1;
    if (src == dest) return 0;

    vector<vector<bool>> vis(n, vector<bool>(m, false));
    queue<tuple<int,int,int>> q;             // {row, col, distance}
    q.push({src.first, src.second, 0});
    vis[src.first][src.second] = true;
    int dr[] = {-1, 1, 0, 0}, dc[] = {0, 0, -1, 1};

    while (!q.empty()) {
        auto [r, c, d] = q.front(); q.pop();
        for (int k = 0; k < 4; k++) {
            int nr = r + dr[k], nc = c + dc[k];
            if (nr < 0 || nr >= n || nc < 0 || nc >= m) continue;
            if (grid[nr][nc] == 0 || vis[nr][nc]) continue;
            if (nr == dest.first && nc == dest.second) return d + 1;
            vis[nr][nc] = true;
            q.push({nr, nc, d + 1});
        }
    }
    return -1;
}

int main() {
    vector<vector<int>> grid = {
        {1, 1, 1, 1},
        {1, 0, 0, 1},
        {1, 1, 1, 1}
    };
    cout << shortestPathBinaryMaze(grid, {0,0}, {2,3}) << endl;  // 5
    return 0;
}
\end{lstlisting}

\subsection*{Dry run}

\dryrun{Source $(0,0)$, destination $(2,3)$ on the $3\times4$ grid
above.}

\begin{itemize}
  \item Layer $0$: $(0,0)$. Layer $1$: $(1,0),(0,1)$.
  \item Layer $2$: $(2,0),(0,2)$. Layer $3$: $(2,1),(0,3)$.
  \item Layer $4$: $(2,2),(1,3)$. Layer $5$: expanding $(2,2)$ or
        $(1,3)$ reaches $(2,3)$ -- return $5$.
  \item The blocked cells $(1,1),(1,2)$ force the path around the
        edges.
\end{itemize}

\subsection*{Complexity}

\begin{complexitybox}
\textbf{Time Complexity.} $O(n\cdot m)$ -- each cell is enqueued and
processed at most once. \textbf{Space Complexity.} $O(n\cdot m)$ for the
visited grid and the BFS queue.
\end{complexitybox}

\begin{takeawaybox}
When every move on a grid costs the same, BFS -- not Dijkstra's -- is the
right tool, and the grid is simply an implicit unit-weight graph. Save
Dijkstra's for grids where cells carry differing move costs, such as
Path With Minimum Effort (Section~\ref{sec:graph-path-min-effort}).
\end{takeawaybox}

\section{Path With Minimum Effort}
\label{sec:graph-path-min-effort}

\problemstatement{Given a 2D grid of heights, find a path from the
top-left cell to the bottom-right cell (moving $4$-directionally)
that minimizes the \textbf{effort} of the path, where a path's effort
is defined as the \textbf{maximum} absolute height difference between
any two consecutive cells along it.}

\begin{patternbox}
\emph{``Minimize the maximum edge weight along a path''} (a
\textbf{minimax path}) is solved with the exact same frontier-expansion
machinery as Section~\ref{sec:graph-dijkstra}'s Dijkstra, with one
change: the quantity tracked and relaxed is no longer a
\textbf{sum} of edge weights, but the \textbf{running maximum} edge
weight seen so far along the path.
\end{patternbox}

\subsection*{Adapting the relaxation rule}

Let $\textit{effort}[\textit{cell}]$ be the minimum possible
``maximum height difference'' achievable on some path from the start
to $\textit{cell}$. Moving from $u$ to a neighbor $v$, the effort of
extending that path is $\max(\textit{effort}[u],
|\textit{height}[v]-\textit{height}[u]|)$ -- not
$\textit{effort}[u] + |\textit{height}[v]-\textit{height}[u]|$, since a
path's effort is its single worst step, not an accumulated total. The
relaxation check becomes: if this candidate effort is
\textbf{smaller} than the currently known $\textit{effort}[v]$, update
and push.

\subsection*{Why Dijkstra's greedy-frontier property still applies}

The same correctness argument from Section~\ref{sec:graph-dijkstra}
carries over unchanged: popping the cell with the smallest known
effort from the priority queue guarantees that effort is final,
because every alternative path to it would have to pass through a
cell with effort $\ge$ the popped value, and taking the
$\max$ with any additional step can only keep the running effort the
same or increase it -- never decrease it. This ``can only stay the
same or grow'' property is the minimax analogue of Dijkstra's
``non-negative weights only add'' requirement.

\subsection*{C++ implementation}

\begin{lstlisting}
#include <bits/stdc++.h>
using namespace std;

int minimumEffortPath(vector<vector<int>>& heights) {
    int rows = heights.size(), cols = heights[0].size();
    vector<vector<int>> effort(rows, vector<int>(cols, INT_MAX));
    effort[0][0] = 0;

    priority_queue<tuple<int,int,int>, vector<tuple<int,int,int>>, greater<>> pq;
    pq.push({0, 0, 0}); // {effort, row, col}

    int dr[] = {-1, 1, 0, 0};
    int dc[] = {0, 0, -1, 1};

    while (!pq.empty()) {
        auto [e, r, c] = pq.top();
        pq.pop();
        if (e > effort[r][c]) continue;
        if (r == rows - 1 && c == cols - 1) return e;

        for (int dir = 0; dir < 4; dir++) {
            int nr = r + dr[dir], nc = c + dc[dir];
            if (nr >= 0 && nr < rows && nc >= 0 && nc < cols) {
                int candidate = max(e, abs(heights[nr][nc] - heights[r][c]));
                if (candidate < effort[nr][nc]) {
                    effort[nr][nc] = candidate;
                    pq.push({candidate, nr, nc});
                }
            }
        }
    }
    return 0;
}

int main() {
    vector<vector<int>> heights = {
        {1, 2, 2},
        {3, 8, 2},
        {5, 3, 5}
    };
    cout << minimumEffortPath(heights) << endl; // 2
    return 0;
}
\end{lstlisting}

\subsection*{Dry run}

\dryrun{Take the $3\times3$ heights grid above.}

The path $(0,0)\to(0,1)\to(0,2)\to(1,2)\to(2,2)$ has consecutive
differences $|2{-}1|=1$, $|2{-}2|=0$, $|2{-}2|=0$, $|5{-}2|=3$ --
effort $3$. A better path,
$(0,0)\to(1,0)\to(2,0)\to(2,1)\to(2,2)$, has differences
$|3{-}1|=2$, $|5{-}3|=2$, $|3{-}5|=2$, $|5{-}3|=2$ -- effort $2$,
worse on individual steps but capped lower overall. Dijkstra's
frontier expansion explores cells in increasing order of their best
known effort, popping $(2,2)$ (the destination) only once its true
minimum effort, $2$, has been finalized.

\begin{complexitybox}
\textbf{Time Complexity.} $O(\textit{rows} \cdot \textit{cols} \cdot
\log(\textit{rows}\cdot\textit{cols}))$, identical structure to
Dijkstra over a grid graph with $4$ edges per node.
\textbf{Space Complexity.} $O(\textit{rows}\cdot\textit{cols})$ for
the \textit{effort} grid and the priority queue.
\end{complexitybox}

\begin{takeawaybox}
Dijkstra's algorithm generalizes to any relaxation rule that is
\textbf{monotonic} -- extending a path by one more edge can only keep
the tracked quantity the same or make it worse, never better. Sum of
weights and running maximum both satisfy this, which is why the exact
same priority-queue frontier expansion solves both the standard
shortest-path problem and this minimax variant.
\end{takeawaybox}

\section{Cheapest Flights Within K Stops}
\label{sec:graph-cheapest-flights-k-stops}

\problemstatement{Given $n$ cities, a list of flights $[from, to,
price]$, a source, a destination, and an integer $k$, find the
cheapest price to travel from source to destination using \textbf{at
most $k$ stops} (i.e.\ at most $k{+}1$ edges). Return $-1$ if no such
route exists.}

\begin{patternbox}
\emph{``Shortest path, but bounded by the number of edges used''}
specifically \textbf{breaks} Dijkstra's greedy frontier-expansion
property: Dijkstra always finalizes whichever node has the globally
smallest distance, with no regard for how many edges that path used --
so it can permanently lock in a cheap-but-too-many-stops path, missing
a more-expensive-but-within-budget alternative. The fix: relax
\textbf{every} edge, for exactly $k{+}1$ \textbf{rounds}, copying
distances between rounds so each round adds at most one new edge to
every path -- this is a bounded version of Bellman-Ford's relaxation
strategy (Section~\ref{sec:graph-bellman-ford}).
\end{patternbox}

\subsection*{Why the copy between rounds is essential}

If edges were relaxed directly into the live $\textit{dist}$ array
within a single round, a single round could chain multiple
relaxations together (relax $u\to v$, then immediately use the updated
$\textit{dist}[v]$ to relax $v \to w$ in the \emph{same} round) --
silently allowing paths with more edges than the round count intends.
Using a fresh copy ($\textit{temp}$) as the round's write target, while
reading from the previous round's frozen $\textit{dist}$, guarantees
each round contributes \textbf{exactly one} additional edge to every
path under consideration.

\subsection*{Algorithm}

\begin{enumerate}
  \item Initialize $\textit{dist}[\textit{src}]=0$, all others
        $\infty$.
  \item Repeat $k{+}1$ times: copy $\textit{dist}$ into
        $\textit{temp}$; for every flight $(u,v,w)$, if
        $\textit{dist}[u] + w < \textit{temp}[v]$, update
        $\textit{temp}[v]$; then set $\textit{dist} \gets \textit{temp}$.
  \item Return $\textit{dist}[\textit{dst}]$ if finite, else $-1$.
\end{enumerate}

\subsection*{C++ implementation}

\begin{lstlisting}
#include <bits/stdc++.h>
using namespace std;

int findCheapestPrice(int n, vector<vector<int>>& flights, int src, int dst, int k) {
    vector<int> dist(n, INT_MAX);
    dist[src] = 0;

    for (int round = 0; round <= k; round++) {
        vector<int> temp = dist;
        for (auto& f : flights) {
            int u = f[0], v = f[1], w = f[2];
            if (dist[u] != INT_MAX && dist[u] + w < temp[v]) {
                temp[v] = dist[u] + w;
            }
        }
        dist = temp;
    }
    return dist[dst] == INT_MAX ? -1 : dist[dst];
}

int main() {
    vector<vector<int>> flights = {{0,1,100},{1,2,100},{0,2,500}};
    cout << findCheapestPrice(3, flights, 0, 2, 1) << endl; // 200
    return 0;
}
\end{lstlisting}

\subsection*{Dry run}

\dryrun{Take $\textit{flights}=[[0,1,100],[1,2,100],[0,2,500]]$,
$\textit{src}=0$, $\textit{dst}=2$, $k=1$ (at most $1$ stop, i.e.\ at
most $2$ edges).}

$\textit{dist}=[0,\infty,\infty]$. Round $0$ ($k{+}1=2$ rounds total):
$\textit{temp}=[0,\infty,\infty]$; relax $0{\to}1$:
$0+100=100<\infty$, $\textit{temp}[1]=100$; relax $1{\to}2$:
$\textit{dist}[1]=\infty$ still (reading the \emph{frozen} pre-round
value), skip; relax $0{\to}2$: $0+500=500<\infty$,
$\textit{temp}[2]=500$. $\textit{dist}=[0,100,500]$. Round $1$:
$\textit{temp}=[0,100,500]$; relax $0{\to}1$: $100$ not $<100$, no
change; relax $1{\to}2$: $\textit{dist}[1]+100=100+100=200<500$,
$\textit{temp}[2]=200$; relax $0{\to}2$: $500$ not $<200$, no change.
$\textit{dist}=[0,100,200]$. Final answer: $\textit{dist}[2]=200$ --
the path $0{\to}1{\to}2$ (one stop, total $200$), correctly preferred
over the direct $0{\to}2$ edge costing $500$.

\begin{complexitybox}
\textbf{Time Complexity.} $O(k \cdot E)$ -- $k{+}1$ rounds, each
scanning every flight once.
\textbf{Space Complexity.} $O(V)$ for \textit{dist} and \textit{temp}.
\end{complexitybox}

\begin{takeawaybox}
This problem is a useful warning against reaching for Dijkstra
reflexively on every ``cheapest path'' phrasing: the moment a
constraint couples the answer to \emph{how the path is structured}
(here, edge count) rather than purely to its total weight, Dijkstra's
greedy finalization can be provably wrong, and a bounded, round-based
relaxation (Bellman-Ford's core idea, capped early) becomes the correct
tool.
\end{takeawaybox}

\section{Number of Ways to Arrive at Destination}
\label{sec:graph-number-of-ways-arrive}

\problemstatement{Given $n$ cities and a list of bidirectional roads
with travel times, count the number of \textbf{distinct shortest-time}
ways to travel from city $0$ to city $n{-}1$, modulo $10^9+7$.}

\begin{patternbox}
\emph{``Count the number of ways to achieve the shortest distance''}
pairs Section~\ref{sec:graph-dijkstra}'s \textit{dist} array with a
parallel \textit{ways} array, updated by the identical
reset-on-improvement, accumulate-on-tie rule the DP on LIS chapter's
Number of LIS used for counting longest chains -- here applied to
Dijkstra's relaxation instead of an array-indexed DP recurrence.
\end{patternbox}

\subsection*{The counting rule, adapted to relaxation}

When relaxing edge $u \to v$ with weight $w$, alongside the usual
distance check:
\begin{itemize}
  \item If $\textit{dist}[u] + w < \textit{dist}[v]$: a
        \textbf{strictly shorter} path to $v$ was just found through
        $u$. Update $\textit{dist}[v]$, and \textbf{reset}
        $\textit{ways}[v] \gets \textit{ways}[u]$ (every previously
        counted way to $v$ used a now-suboptimal distance).
  \item If $\textit{dist}[u] + w = \textit{dist}[v]$: this is
        \textbf{another equally short} path to $v$. \textbf{Accumulate}:
        $\textit{ways}[v] \mathrel{+}= \textit{ways}[u]$.
\end{itemize}

\subsection*{C++ implementation}

\begin{lstlisting}
#include <bits/stdc++.h>
using namespace std;
const int MOD = 1e9 + 7;

int countPaths(int n, vector<vector<int>>& roads) {
    vector<vector<pair<long long,long long>>> adj(n); // {neighbor, weight}
    for (auto& r : roads) {
        adj[r[0]].push_back({r[1], r[2]});
        adj[r[1]].push_back({r[0], r[2]});
    }

    vector<long long> dist(n, LLONG_MAX), ways(n, 0);
    dist[0] = 0;
    ways[0] = 1;
    priority_queue<pair<long long,int>, vector<pair<long long,int>>, greater<>> pq;
    pq.push({0, 0});

    while (!pq.empty()) {
        auto [d, node] = pq.top();
        pq.pop();
        if (d > dist[node]) continue;

        for (auto& [neighbor, weight] : adj[node]) {
            if (dist[node] + weight < dist[neighbor]) {
                dist[neighbor] = dist[node] + weight;
                ways[neighbor] = ways[node];
                pq.push({dist[neighbor], neighbor});
            } else if (dist[node] + weight == dist[neighbor]) {
                ways[neighbor] = (ways[neighbor] + ways[node]) % MOD;
            }
        }
    }
    return (int)(ways[n - 1] % MOD);
}

int main() {
    vector<vector<int>> roads = {{0,6,7},{0,1,2},{1,2,3},{1,3,3},
                                  {6,3,3},{3,5,1},{6,5,1},{2,5,1},{0,4,5},{4,6,2}};
    cout << countPaths(7, roads) << endl; // 4
    return 0;
}
\end{lstlisting}

\subsection*{Dry run}

\dryrun{Take a small graph: $0{-}1$ (weight $1$), $0{-}2$ (weight
$1$), $1{-}3$ (weight $1$), $2{-}3$ (weight $1$) -- two equally short
paths from $0$ to $3$.}

$\textit{dist}=[0,\infty,\infty,\infty]$, $\textit{ways}=[1,0,0,0]$.
Pop $0$: relax $1$ ($\infty\to1$, reset $\textit{ways}[1]=1$); relax
$2$ ($\infty\to1$, reset $\textit{ways}[2]=1$). Pop $1$ (dist $1$):
relax $3$ ($\infty\to2$, reset $\textit{ways}[3]=\textit{ways}[1]=1$).
Pop $2$ (dist $1$): relax $3$ -- $\textit{dist}[2]+1=2=\textit{dist}[3]$,
\textbf{tie}, accumulate $\textit{ways}[3] \mathrel{+}=
\textit{ways}[2]=1$, giving $\textit{ways}[3]=2$. Pop $3$: no further
relaxation. Final: $2$ shortest paths to node $3$ (via $1$, and via
$2$), correctly counted.

\begin{complexitybox}
\textbf{Time Complexity.} $O((V+E)\log V)$, identical to plain
Dijkstra -- the \textit{ways} bookkeeping adds only $O(1)$ work per
relaxation.
\textbf{Space Complexity.} $O(V+E)$ for the adjacency list, plus
$O(V)$ for \textit{dist} and \textit{ways}.
\end{complexitybox}

\begin{takeawaybox}
The ``exists / count / optimize'' framing introduced in the DP on
Subsequences chapter, and reused for counting longest chains in the DP
on LIS chapter, extends naturally beyond dynamic programming tables --
any algorithm computing an optimum via relaxation or comparison (here,
Dijkstra's distance updates) can be paired with a parallel counting
array using the exact same reset-on-improvement,
accumulate-on-tie rule.
\end{takeawaybox}

\section{Minimum Multiplications to Reach End}
\label{sec:graph-min-multiplications}

\problemstatement{Given a start number, an end number, and an array of
multipliers, in each step you may replace the current number $x$ with
$(x \times \textit{arr}[i]) \bmod 100000$ for any $i$. Find the minimum
number of steps to turn \textit{start} into \textit{end}, or $-1$ if
impossible.}

\begin{patternbox}
The disguise here is strong: there is no explicit graph. But each
reachable number $0\ldots99999$ is a \textbf{node}, and each multiplier
defines an \textbf{edge} of cost one step. ``Minimum steps, all costs
equal'' is again a \textbf{unit-weight shortest path}, so \textbf{BFS}
over the $100000$ possible numbers solves it.
\end{patternbox}

\subsection*{Modeling numbers as nodes}

The modulo keeps every reachable value in $[0,99999]$, a finite state
space. From a value $x$, applying multiplier $\textit{arr}[i]$ produces
a neighbor $(x\cdot\textit{arr}[i])\bmod 100000$. Every such transition
is one step, so all edges have equal weight and BFS from \textit{start}
finds the fewest steps to \textit{end}. A visited array over the
$100000$ states prevents revisiting and guarantees termination.

\begin{ideabox}[A Bounded State Space Makes BFS Finite]
The $\bmod\,100000$ is what makes this tractable: without it the numbers
could grow without bound. With it, there are only $100000$ possible
states, so BFS explores each at most once and either reaches
\textit{end} or exhausts the reachable set.
\end{ideabox}

\subsection*{Algorithm}

\begin{enumerate}
  \item If \textit{start} equals \textit{end}, return $0$. Push
        \textit{start} with $\textit{steps}=0$; mark it visited.
  \item While the queue is non-empty, pop $(\textit{value},
        \textit{steps})$.
  \item For each multiplier, compute
        $\textit{next}=(\textit{value}\cdot\textit{arr}[i])\bmod
        100000$; if $\textit{next}=\textit{end}$, return
        $\textit{steps}+1$; if unvisited, mark and push it with
        $\textit{steps}+1$.
  \item If the queue empties, return $-1$.
\end{enumerate}

\subsection*{C++ implementation}

\begin{lstlisting}
#include <bits/stdc++.h>
using namespace std;

int minimumMultiplications(vector<int>& arr, int start, int end) {
    if (start == end) return 0;
    const int MOD = 100000;
    vector<bool> vis(MOD, false);
    queue<pair<int,int>> q;                  // {value, steps}
    q.push({start, 0});
    vis[start] = true;

    while (!q.empty()) {
        auto [value, steps] = q.front(); q.pop();
        for (int m : arr) {
            int next = (int)((long long)value * m % MOD);
            if (next == end) return steps + 1;
            if (!vis[next]) {
                vis[next] = true;
                q.push({next, steps + 1});
            }
        }
    }
    return -1;
}

int main() {
    vector<int> arr = {2, 5, 7};
    cout << minimumMultiplications(arr, 3, 30) << endl;   // 2  (3*2=6, 6*5=30)
    return 0;
}
\end{lstlisting}

\subsection*{Dry run}

\dryrun{\textit{arr}$=\{2,5,7\}$, start $3$, end $30$.}

\begin{itemize}
  \item Step $0$: value $3$. Neighbors: $6$ (via $2$), $15$ (via $5$),
        $21$ (via $7$) -- none is $30$; enqueue all with steps $1$.
  \item Step $1$, value $6$: $6\times2=12$, $6\times5=30$ --
        \textbf{match!} Return $1+1=2$.
  \item Shortest: $3\xrightarrow{\times2}6\xrightarrow{\times5}30$, two
        steps.
\end{itemize}

\subsection*{Complexity}

\begin{complexitybox}
\textbf{Time Complexity.} $O(100000 \times \lvert\textit{arr}\rvert)$ --
each of the $100000$ states is expanded once, trying every multiplier.
\textbf{Space Complexity.} $O(100000)$ for the visited array and queue.
\end{complexitybox}

\begin{takeawaybox}
When a problem asks for the minimum number of equal-cost transformations
between two values, model each value as a node and each transformation
as a unit edge, then run BFS -- even when no graph is stated. A bounded
state space (here, the modulo) is the signal that this modeling is
finite and BFS applies.
\end{takeawaybox}

\section{Bellman-Ford Algorithm}
\label{sec:graph-bellman-ford}

\problemstatement{Given a weighted directed graph, possibly containing
\textbf{negative} edge weights, find the shortest distance from a
source to every other node -- or detect that a \textbf{negative-weight
cycle} makes the notion of ``shortest path'' undefined.}

\begin{patternbox}
\emph{``Shortest path, negative weights allowed''} is solved by
relaxing \textbf{every edge}, $V{-}1$ times over (no priority queue, no
greedy frontier choice -- Dijkstra's greedy finalization is provably
unsafe the moment a negative edge exists, as discussed in
Section~\ref{sec:graph-dijkstra}). A final, $V$-th round that still
finds an improvement reveals a negative cycle.
\end{patternbox}

\subsection*{Why $V-1$ rounds are exactly enough}

In a graph with no negative cycle, the shortest path between any two
nodes is always a \textbf{simple} path (visits no node twice) -- a path
revisiting a node would include a cycle, and removing that cycle could
only shorten the path further unless the cycle's total weight is
\emph{exactly} $0$, in which case removing it doesn't change the
length either way. A simple path in a graph with $V$ nodes has at most
$V{-}1$ edges. Each full round of relaxing every edge extends the
\emph{confirmed-correct} prefix of every shortest path by (at least)
one more edge -- so after $V{-}1$ rounds, every shortest path (having
at most $V{-}1$ edges) is guaranteed fully relaxed and correct.

\subsection*{Why a $V$-th round detects a negative cycle}

If, after $V{-}1$ rounds (when every legitimate shortest path must
already be finalized), a $V$-th round of relaxation \emph{still} finds
an improvement on some edge, that improvement cannot correspond to any
genuine shortest path (those are already settled) -- it can only be
explained by a cycle whose total weight is negative, allowing
``distance'' to be driven arbitrarily low by looping through it
repeatedly.

\subsection*{Algorithm}

\begin{enumerate}
  \item Initialize $\textit{dist}[\textit{src}]=0$, all others
        $\infty$.
  \item Repeat $V{-}1$ times: for every edge $(u,v,w)$, if
        $\textit{dist}[u] \neq \infty$ and $\textit{dist}[u]+w <
        \textit{dist}[v]$, update $\textit{dist}[v]$.
  \item Run one more relaxation pass over all edges; if any edge can
        still be relaxed, report a negative cycle.
\end{enumerate}

\subsection*{C++ implementation}

\begin{lstlisting}
#include <bits/stdc++.h>
using namespace std;

vector<int> bellmanFord(int n, vector<vector<int>>& edges, int src) {
    vector<int> dist(n, INT_MAX);
    dist[src] = 0;

    for (int i = 0; i < n - 1; i++) {
        for (auto& e : edges) {
            int u = e[0], v = e[1], w = e[2];
            if (dist[u] != INT_MAX && dist[u] + w < dist[v]) {
                dist[v] = dist[u] + w;
            }
        }
    }

    for (auto& e : edges) {
        int u = e[0], v = e[1], w = e[2];
        if (dist[u] != INT_MAX && dist[u] + w < dist[v]) {
            return {-1}; // sentinel: negative cycle detected
        }
    }
    return dist;
}

int main() {
    vector<vector<int>> edges = {{0,1,4},{0,2,5},{1,2,-3},{2,3,4}};
    vector<int> dist = bellmanFord(4, edges, 0);
    for (int d : dist) cout << d << " ";
    cout << endl; // 0 4 1 5
    return 0;
}
\end{lstlisting}

\subsection*{Dry run}

\dryrun{Take edges $0{\to}1(4)$, $0{\to}2(5)$, $1{\to}2(-3)$,
$2{\to}3(4)$, $n=4$, source $0$.}

$\textit{dist}=[0,\infty,\infty,\infty]$. Round $1$: relax $0{\to}1$
($\infty\to4$); relax $0{\to}2$ ($\infty\to5$); relax $1{\to}2$
($4{-}3=1<5$, update to $1$); relax $2{\to}3$ ($1+4=5<\infty$, update
to $5$). $\textit{dist}=[0,4,1,5]$. Round $2$: relax $0{\to}1$ ($4$
not $<4$); relax $0{\to}2$ ($5$ not $<1$); relax $1{\to}2$ ($1$ not
$<1$); relax $2{\to}3$ ($5$ not $<5$) -- no changes,
\textbf{stabilized} early (rounds can stop as soon as a full pass makes
no change, though the algorithm above runs the full $V{-}1$ rounds
regardless). Round $3$ ($n{-}1=3$ total): no changes. Final check
round: no edge relaxes further -- no negative cycle. Final:
$[0,4,1,5]$.

\begin{complexitybox}
\textbf{Time Complexity.} $O(V \cdot E)$ -- $V{-}1$ rounds (plus one
check round), each scanning all $E$ edges.
\textbf{Space Complexity.} $O(V)$ for \textit{dist}.
\end{complexitybox}

\begin{takeawaybox}
Bellman-Ford trades Dijkstra's $O((V+E)\log V)$ speed for two things
Dijkstra cannot offer: correctness with negative edges, and the ability
to detect negative cycles at all. When a problem's weights are
guaranteed non-negative, Dijkstra remains the faster default --
Bellman-Ford is the fallback specifically for when that guarantee is
unavailable.
\end{takeawaybox}

\section{Floyd-Warshall Algorithm}
\label{sec:graph-floyd-warshall}

\problemstatement{Given a weighted directed graph (as an
$n \times n$ adjacency matrix, with $\infty$ for missing edges), find
the shortest distance between \textbf{every pair} of nodes.}

\begin{patternbox}
\emph{``Shortest path between every pair of nodes, simultaneously''}
is solved not by running Section~\ref{sec:graph-dijkstra}'s Dijkstra
from every node (which would cost $O(V \cdot (V+E)\log V)$), but by a
direct dynamic programming recurrence over which nodes are
\textbf{allowed as intermediate stops}: progressively allow node $1$,
then nodes $1$ and $2$, then $1,2,3$, and so on, updating every pair's
shortest distance each time a new node becomes available as a
potential stop along the way.
\end{patternbox}

\subsection*{Deriving the recurrence}

Let $\textit{dist}_k[i][j]$ be the shortest distance from $i$ to $j$
using \textbf{only} nodes $1,\ldots,k$ as intermediate stops (besides
$i,j$ themselves). Either the shortest such path doesn't use node $k$
at all (so $\textit{dist}_k[i][j] = \textit{dist}_{k-1}[i][j]$), or it
does use $k$ -- exactly once, as a single stop, splitting the path into
an $i\to k$ segment and a $k\to j$ segment, both of which only need
nodes $1,\ldots,k{-}1$ as their own intermediates (since they don't
revisit $k$):
\[
  \textit{dist}_k[i][j] = \min\big(\textit{dist}_{k-1}[i][j],\
  \textit{dist}_{k-1}[i][k] + \textit{dist}_{k-1}[k][j]\big).
\]
Implemented in place, with $k$ as the \textbf{outermost} loop, this
needs only a single 2D array (each cell is read and written using
already-updated values from the current and prior $k$, which is safe
because $\textit{dist}[i][k]$ and $\textit{dist}[k][j]$ are never
themselves updated using $k$ as an intermediate -- a path from $i$ to
$k$ using $k$ itself as a stop would be nonsensical).

\subsection*{C++ implementation}

\begin{lstlisting}
#include <bits/stdc++.h>
using namespace std;
const int INF = 1e9;

void floydWarshall(vector<vector<int>>& dist) {
    int n = dist.size();

    for (int k = 0; k < n; k++) {
        for (int i = 0; i < n; i++) {
            for (int j = 0; j < n; j++) {
                if (dist[i][k] < INF && dist[k][j] < INF) {
                    dist[i][j] = min(dist[i][j], dist[i][k] + dist[k][j]);
                }
            }
        }
    }
}

int main() {
    int n = 4;
    vector<vector<int>> dist(n, vector<int>(n, INF));
    for (int i = 0; i < n; i++) dist[i][i] = 0;
    dist[0][1] = 5; dist[1][2] = 3; dist[2][3] = 1; dist[0][3] = 100;

    floydWarshall(dist);
    cout << dist[0][3] << endl; // 9  (via 0->1->2->3: 5+3+1)
    return 0;
}
\end{lstlisting}

\subsection*{Dry run}

\dryrun{Take the $4$-node graph above: $0{\to}1(5)$, $1{\to}2(3)$,
$2{\to}3(1)$, $0{\to}3(100)$.}

Initial $\textit{dist}[0][3]=100$ (direct edge). After allowing $k=1$:
$\textit{dist}[0][2] = \min(\infty, \textit{dist}[0][1]+\textit{dist}[1][2]) =
\min(\infty,5+3)=8$. After allowing $k=2$: $\textit{dist}[0][3] =
\min(100, \textit{dist}[0][2]+\textit{dist}[2][3]) = \min(100,8+1)=9$.
After allowing $k=3$: no further improvement possible through $3$
(it has no outgoing edges in this example). Final
$\textit{dist}[0][3]=9$, matching the path $0\to1\to2\to3$.

\begin{complexitybox}
\textbf{Time Complexity.} $O(V^3)$ -- three nested loops over all
nodes.
\textbf{Space Complexity.} $O(V^2)$ for the distance matrix (updated
in place; no extra space needed per layer of $k$).
\end{complexitybox}

\begin{takeawaybox}
Floyd-Warshall is this book's clearest example of an
\textbf{ending-here}-style DP (Section~\ref{sec:graph-shortest-path-intro}'s
opening idea, generalized) applied to a 2D state $(i,j)$ with a third
dimension -- ``which nodes are allowed as intermediates'' -- collapsed
into a single in-place array by processing $k$ in the outer loop, the
same in-place-update trick used by every space-optimized DP table
throughout this book series.
\end{takeawaybox}

\section{City With the Smallest Number of Neighbors at a Threshold Distance}
\label{sec:graph-city-smallest-neighbors}

\problemstatement{Given $n$ cities connected by weighted edges and a
$\textit{distanceThreshold}$, find the city from which the
\textbf{fewest} other cities are reachable within
$\textit{distanceThreshold}$. If multiple cities tie for fewest,
return the one with the \textbf{largest} index.}

\begin{patternbox}
This closing problem of the chapter is a direct application of
Section~\ref{sec:graph-floyd-warshall}: once \textbf{all-pairs}
shortest distances are known, answering ``how many cities are within
threshold $T$ of city $i$'' for every $i$ is a simple linear scan over
each row of the distance matrix -- no further graph algorithm is
needed.
\end{patternbox}

\subsection*{Reduction}

\begin{enumerate}
  \item Build the $n \times n$ distance matrix and run Floyd-Warshall
        to fill in all-pairs shortest distances.
  \item For each city $i$, count how many cities $j \neq i$ satisfy
        $\textit{dist}[i][j] \le \textit{distanceThreshold}$.
  \item Track the city with the smallest count, breaking ties by
        preferring the \textbf{larger} index (scan $i$ from $0$ to
        $n{-}1$, and use $\le$ rather than $<$ when comparing counts,
        so a later equal-or-better city always overwrites an earlier
        one).
\end{enumerate}

\subsection*{C++ implementation}

\begin{lstlisting}
#include <bits/stdc++.h>
using namespace std;
const int INF = 1e9;

int findTheCity(int n, vector<vector<int>>& edges, int distanceThreshold) {
    vector<vector<int>> dist(n, vector<int>(n, INF));
    for (int i = 0; i < n; i++) dist[i][i] = 0;
    for (auto& e : edges) {
        dist[e[0]][e[1]] = e[2];
        dist[e[1]][e[0]] = e[2];
    }

    for (int k = 0; k < n; k++) {
        for (int i = 0; i < n; i++) {
            for (int j = 0; j < n; j++) {
                if (dist[i][k] < INF && dist[k][j] < INF) {
                    dist[i][j] = min(dist[i][j], dist[i][k] + dist[k][j]);
                }
            }
        }
    }

    int bestCity = -1, minCount = n;
    for (int i = 0; i < n; i++) {
        int count = 0;
        for (int j = 0; j < n; j++) {
            if (i != j && dist[i][j] <= distanceThreshold) count++;
        }
        if (count <= minCount) { // <= prefers the larger index on ties
            minCount = count;
            bestCity = i;
        }
    }
    return bestCity;
}

int main() {
    vector<vector<int>> edges = {{0,1,3},{1,2,1},{1,3,4},{2,3,1}};
    cout << findTheCity(4, edges, 4) << endl; // 3
    return 0;
}
\end{lstlisting}

\subsection*{Dry run}

\dryrun{Take $n=4$, edges $0{-}1(3),1{-}2(1),1{-}3(4),2{-}3(1)$,
$\textit{distanceThreshold}=4$.}

All-pairs distances after Floyd-Warshall: from $0$: $[0,3,4,5]$ (to
$2$: $3+1=4$ via $1$; to $3$: $3+1+1=5$ via $1,2$, shorter than the
direct $1{-}3$ edge route $3+4=7$). From $1$: $[3,0,1,2]$. From $2$:
$[4,1,0,1]$. From $3$: $[5,2,1,0]$ (to $1$: $1+1=2$ via $2$, shorter
than the direct edge $4$). Counts within threshold $4$ (excluding
self): city $0$ sees $1$ (dist $3$) and $2$ (dist $4$) -- $2$ cities
($3$ at distance $5$ is excluded). City $1$ sees all three others --
$3$ cities. City $2$ sees all three others -- $3$ cities. City $3$
sees $1$ (dist $2$) and $2$ (dist $1$) -- $2$ cities ($0$ at distance
$5$ is excluded). Cities $0$ and $3$ tie at $2$ neighbors each; the
tie-break (larger index) selects city $3$.

\begin{complexitybox}
\textbf{Time Complexity.} $O(V^3)$ for Floyd-Warshall, plus $O(V^2)$
for the counting scan: $O(V^3)$ overall.
\textbf{Space Complexity.} $O(V^2)$ for the distance matrix.
\end{complexitybox}

\begin{takeawaybox}
This chapter closes by demonstrating exactly why Floyd-Warshall earns
its $O(V^3)$ cost despite being asymptotically worse than running
Dijkstra from every node ($O(V(V+E)\log V)$ for sparse graphs): once
\emph{every} pairwise distance is available at once, follow-up
questions like this one (counting neighbors within a threshold, for
every city) become trivial linear scans -- exactly the same trade-off
this book has made repeatedly between an upfront, fuller computation
and cheap, flexible queries afterward.
\end{takeawaybox}

